\section{SAT Method}
\label{sec:sat-method}

\subsection{Decision formulation}
Let $G=(V,E)$ be a simple undirected graph, with $n=|V|$ and $m=|E|$. A
labeling is a bijection $f:V\rightarrow\{1,\ldots,n\}$. Its antibandwidth is
\begin{equation}
  B(f)=\min_{\{u,v\}\in E}|f(u)-f(v)|,
\end{equation}
and the optimization problem is to maximize $B(f)$. For a fixed threshold
$K\geq 1$, the SAT decision problem asks whether there exists a bijection
satisfying
\begin{equation}
  |f(u)-f(v)|\geq K
  \qquad \forall\,\{u,v\}\in E.
  \label{eq:sat-decision}
\end{equation}
The implementation uses two alternative encodings of this same decision
problem: a pairwise assignment encoding and an order encoding. Both include
constraints for vertex assignment, label permutation, and edge distance.
They differ in their representation of a vertex label.

\subsection{Pairwise assignment encoding}
For each vertex $v\in V$ and label $p\in\{1,\ldots,n\}$, introduce a Boolean
variable $x_{v,p}$, where $x_{v,p}=1$ means that vertex $v$ receives label $p$.
The encoding contains three groups of clauses.

\paragraph{Assignment.}
Each vertex receives at least one label (ALO), and no vertex receives two
different labels (AMO):
\begin{align}
  &\bigvee_{p=1}^{n}x_{v,p}
      &&\forall v\in V,\\
  &\neg x_{v,p}\vee\neg x_{v,q}
      &&\forall v\in V,\;1\leq p<q\leq n.
\end{align}

\paragraph{Permutation.}
Each label is assigned to at least one vertex and to at most one vertex:
\begin{align}
  &\bigvee_{v\in V}x_{v,p}
      &&\forall p\in\{1,\ldots,n\},\\
  &\neg x_{u,p}\vee\neg x_{v,p}
      &&\forall p\in\{1,\ldots,n\},\;u<v.
\end{align}
Together, assignment and permutation make the labeling a bijection.

\paragraph{Edge distance.}
For every edge $\{u,v\}\in E$, every label pair $(p,q)$ whose separation is
strictly smaller than $K$ is forbidden:
\begin{equation}
  \neg x_{u,p}\vee\neg x_{v,q}
  \qquad
  \forall\{u,v\}\in E,\quad
  1\leq |p-q|\leq K-1.
  \label{eq:pairwise-distance}
\end{equation}
For $1\leq K<n$, the number of edge-distance clauses is
\begin{equation}
  2m\sum_{d=1}^{K-1}(n-d).
\end{equation}
With pairwise AMO, the assignment and permutation groups each contain
$n+n\binom{n}{2}$ clauses. The pairwise encoding uses $n^2$ primary variables;
an alternative AMO encoding may introduce auxiliary variables.

\subsection{Order encoding}
The order encoding uses Boolean variables $o_{v,p}$ for
$v\in V$ and $p\in\{1,\ldots,n-1\}$, with the meaning
\begin{equation}
  o_{v,p}=1 \quad\Longleftrightarrow\quad f(v)\leq p.
  \label{eq:order-variable}
\end{equation}
The boundary values are implicit: $o_{v,0}=0$ and $o_{v,n}=1$. For every
vertex, monotonicity is enforced by
\begin{equation}
  \neg o_{v,p}\vee o_{v,p+1}
  \qquad
  \forall v\in V,\;p=1,\ldots,n-2.
  \label{eq:order-chain}
\end{equation}
A monotone chain uniquely determines one label: the smallest index $p$ for
which $o_{v,p}=1$, or $n$ if all variables are false.

To express that a vertex does not receive label $q$, define the disjunction
\begin{equation}
  N(v,q)=
  \begin{cases}
    \neg o_{v,1}, & q=1,\\
    \neg o_{v,q}\vee o_{v,q-1}, & 1<q<n,\\
    o_{v,n-1}, & q=n.
  \end{cases}
  \label{eq:not-label}
\end{equation}
For every pair of distinct vertices $u<v$ and every label $q$, the
permutation constraint is
\begin{equation}
  N(u,q)\vee N(v,q).
  \label{eq:order-permutation}
\end{equation}
Thus, two vertices cannot receive the same label. Since there are $n$ vertices
and $n$ available labels, injectivity also implies surjectivity.

The edge-distance constraints use the same no-good construction:
\begin{equation}
  N(u,p)\vee N(v,q)
  \qquad
  \forall\{u,v\}\in E,\quad
  1\leq |p-q|\leq K-1.
  \label{eq:order-distance}
\end{equation}
The order encoding has $n(n-1)$ variables, $n(n-2)$ chain clauses, and
$n\binom{n}{2}$ permutation clauses. Its edge-distance clause count is the
same as in the pairwise encoding, although the clauses can contain up to four
literals. The order encoding does not use a separate AMO constraint: the
monotone chain replaces the ALO/AMO pair. Therefore, changing the AMO metadata
for the order encoding does not change its CNF.

\subsection{AMO variants}
For the pairwise assignment encoding, the default AMO formulation adds
$\binom{r}{2}$ binary clauses for a group of $r$ literals and introduces no
auxiliary variables. The sequential-counter variant instead delegates the
at-most-one constraint to \texttt{pysat.card} and adds auxiliary variables to
reduce the number of clauses. Both variants preserve the same feasible
labelings when the auxiliary variables are existentially eliminated. AMO is
not an experimental factor for the order encoding, because the order chain
already represents one label per vertex.

\subsection{SAT backend and optimization search}
The SAT implementation uses the PySAT interface and supports CaDiCaL
(\texttt{cadical195}) and Glucose (\texttt{glucose4}) as the Protocol Core
solver candidates. Kissat is excluded from the Core after the tested PySAT
binding exhibited incorrect model/assumption behavior and process failures.
A fresh backend is created for each threshold probe.

The optimization procedure solves a sequence of decision problems. It keeps
$\ell$, the best verified antibandwidth found so far, and $h$, a threshold
already known to be infeasible. Initially $\ell=0$ and $h=n$, because no
labeling can have a difference of $n$ or greater. While $\ell+1<h$, the next
threshold is $K=\ell+1$:
\begin{itemize}
  \item If the SAT solver returns SAT, the model is decoded and independently
    verified. If its actual antibandwidth is $B$, the algorithm sets
    $\ell=B$ and starts the next probe at $B+1$.
  \item If the solver returns UNSAT, it sets $h=K$. Since $K=\ell+1$, the
    current verified labeling is optimal.
  \item If the solver returns UNKNOWN or the run budget expires, the search
    stops. A verified incumbent is reported as FEAS; without one, the result is
    TO.
\end{itemize}
A SAT labeling with value $B\geq K$ also satisfies every threshold between
$K$ and $B$. Jumping directly to $B+1$ therefore skips only thresholds already
known to be feasible. The first infeasible threshold, when found, proves
optimality. The initial counting bound also proves optimality if a labeling
reaches $n-1$.

The time budget covers CNF construction, clause loading, all solver probes,
decoding, and verification for the entire K-search sequence. Each probe
receives the remaining budget. The shared runner adds a hard process timeout
and is responsible for enforcing the configured OS-level memory limit.
\texttt{peak\_memory} is a sampled process-memory measurement, not a
continuous high-water mark; the official protocol fixes one Windows host so
that the reported memory measurements are taken under a consistent
environment.

\subsection{Decoding, verification, and result semantics}
A SAT model is not accepted directly as a solution. The encoding-specific
decoder must reconstruct a permutation of $\{1,\ldots,n\}$. The shared
\texttt{verify} routine then independently checks the labeling and recomputes
\[
  B(f)=\min_{\{u,v\}\in E}|f(u)-f(v)|.
\]
The recomputed value must meet the probed threshold and agree with the
reported lower bound. The project uses the following status semantics:
\texttt{OPT} denotes a verified incumbent with an optimality proof;
\texttt{FEAS} denotes a verified incumbent without a proof at termination;
\texttt{TO} denotes termination without a verified incumbent; \texttt{MO}
denotes a confirmed memory-limit event; and \texttt{ERR} denotes an execution,
backend, license, malformed-output, or correctness error. In particular, a
license failure is not reported as a timeout.

\subsection{Correctness evidence}
The W2 correctness matrix evaluated the two ABP encodings, two AMO variants,
and two SAT solvers (eight configurations) against a brute-force oracle on
small graphs. The test suite also checks SAT/UNSAT decisions over thresholds,
model decoding, model counts, and the AMO contract. The empirical matrix is
evidence for the tested finite instances, not a proof that the SAT backends
always answer UNSAT correctly. The proof argument for the encodings and the
K-search invariants is provided in Appendix~\ref{app:sat-correctness}.
